\section{2026 年 11 月核验增补}

本节按论文首次提交时间排列。每项方法附有可核验的外部入选依据。

% METHOD: 2605.25563
\paragraph{CodecSplat：编码规则的二维生成潜特征而非生成后的高斯}
2026.05~\cite{2605_25563}。CodecSplat 将压缩点置于前馈高斯生成网络内部。编码器压缩尚具二维规则排列的中间特征，接收端从重建特征回归高斯位置、几何与外观，可避免对已有不规则原语逐项量化。

作者在 DL3DV 和 RE10K 对比先生成再编码的管线与其率失真表现。训练分为高斯预测与率失真训练两个阶段，基于已发布的 DepthSplat 初始化。低码率结论针对论文的前馈重建协议，不代表可直接压缩任意已经生成的资产。作者公开了源代码及运行脚本。

\textbf{入选依据} \href{https://github.com/pengpeng-yu/codecsplat}{公开实现}已找到与论文对应的代码文件。

% METHOD: 2609.03657
\paragraph{用三维高斯扰动构造视频修复训练对}
2026.09~\cite{2609_03657}。为学习能修复跨视图伪影的视频先验，通常要从多种欠观测条件反复拟合场景，再渲染训练视频。该工作直接扰动已拟合高斯的缩放、旋转和保留状态，生成带空间对应的退化视频，与同场景的干净渲染配对。免优化指构造退化场景的过程，视频先验仍需训练。

论文在多个稀疏观测设置比较二维退化和不同三维扰动，另用扩展模型评价深度与下游策略表现。三维扰动提供的是一种训练数据生成近似，对特定相机、退化分布和重建问题的覆盖仍应在目标数据上验证。

\textbf{入选依据} 论文署名单位包括 Huawei、Tavus（\href{https://arxiv.org/abs/2609.03657}{论文原文}）。

% METHOD: 2610.19502
\paragraph{Holi-Spatial 360：用匹配约束稀疏全景重建并生成空间标注}
2026.10~\cite{2610_19502}。稀疏全景相距较远时，图像修复容易改变房间原有结构。Holi-Spatial 360 用 DA3 的匹配给高斯优化提供几何锚定，再用适配的图像与视频模型补充连续轨迹视图。重建和 RGB-D 验证还被用于生成三维实例、描述与空间问答，构成可用于空间推理的数据。

10 个 RealSee3D 场景的 LiDAR 对照中，各方法采用相同 DA3 初始化，2 厘米阈值的 F-score 从 DA3 的 0.2207 提高到 0.2617。该检验覆盖的是这些室内场景的几何，未观测区域的生成内容仍受对应关系可靠性和参考覆盖限制。

\textbf{入选依据} 论文署名单位包括 Tencent（\href{https://arxiv.org/html/2610.19502v1}{论文原文}）。

% METHOD: 2610.21634
\paragraph{Anchor3DGS：以共享三维锚点替代逐像素生成}
2026.10~\cite{2610_21634}。Anchor3DGS 从未标定输入构造几何支架，以固定数量的三维锚点聚合可见多视图信息，在锚点之间交换局部结构上下文，再把锚点解码为局部高斯。锚点预算不再随像素或视图数直接增长，并通过熵引导将更多容量分配给复杂区域。

论文以 3 和 6 个输入视图在 VR-NeRF、MipNeRF360 与独立的 DL3DV 子集评价，在 MipNeRF360 与 DL3DV 上报告超过 2 dB 的 PSNR 增益。该数字对应论文所列的稀疏设置和基线；极少视图中的结构仍受支架质量、锚点分辨率和遮挡影响。

\textbf{入选依据} 论文署名单位包括 Ant Group（\href{https://arxiv.org/abs/2610.21634}{论文原文}）。

% METHOD: 2611.26431
\paragraph{FFGaussType：把已知字形轮廓转为可复用高斯资产}
2026.11~\cite{2611_26431}。FFGaussType 面向已知轮廓的挤出字形与 SVG 图标。方法从采样锚点预测二维特征平面，通过跨资产共享的解码器把平面特征转成带法向的神经高斯。相比逐资产拟合，学习的对象是从给定轮廓几何到高斯资产的转换。

论文按字体家族或图标来源拆分数据，报告全测试集法向误差 1.62 度，来自已准备锚点的转换耗时为 9.1 毫秒。该时间不含轮廓锚点的前处理。论文仅研究具有均一材质的挤出资产，着色不处理自阴影或与周边场景的光传输。

\textbf{入选依据} 论文署名单位包括 Dolby Laboratories（\href{https://arxiv.org/abs/2611.26431}{论文原文}）。
